\begin{helper}
Let \(G\) be a loop graph and let \(D^*(G)\) be its star-product digraph.
Fix a marking rule on children of forward independent prefixes, and let
\(\operatorname{sig}\) be the marked-tuple signature from
\(\noderef{StarProductMarkedTupleSignature}\).  Let \(k,w,\Delta,h\) be
nonnegative integers with \(w\le k\) and \(h\le\Delta\).  Suppose that every
forward independent \(k\)-tuple \(v\) in \(D^*(G)\) has at most \(w\)
unmarked steps, that is,
\[
  |\operatorname{sig}(v)|\le w.
\]
Assume also that every forward independent prefix of length \(m<k\) has at
most \(\Delta\) legal forward-independent children and at most \(h\) legal
marked children.  Then
\[
  \fwi_k(D^*(G))\le 2^k\Delta^w h^{k-w}.
\]
\end{helper}
\begin{proof}
Let \(P\) be the finite set of forward independent \(k\)-tuples in
\(D^*(G)\).  By \(\noderef{ForwardIndependentTupleCount}\), its cardinality
is \(\fwi_k(D^*(G))\).  Define
\[
  \operatorname{sig}:P\to\{0,1\}^k
\]
by the marked-tuple signature of \(\noderef{StarProductMarkedTupleSignature}\).
The assumed path condition says exactly that every \(p\in P\) has signature
weight at most \(w\).

It remains to verify the fiber hypothesis of
\(\noderef{MarkedTreePathCounting}\).  Fix a binary string \(z\).  The fiber
\(\{p\in P:\operatorname{sig}(p)=z\}\) is precisely the fixed-signature
star-product fiber counted by \(\noderef{StarProductMarkedTupleFiberCount}\).
The prefix hypotheses on all legal children and on marked legal children are
therefore exactly the hypotheses of
\(\noderef{StarProductMarkedTupleFiberBound}\), which gives
\[
  |\{p\in P:\operatorname{sig}(p)=z\}|
  \le \Delta^{|z|}h^{k-|z|}.
\]

Now all hypotheses of \(\noderef{MarkedTreePathCounting}\) hold: \(w\le k\),
\(h\le\Delta\), every path has at most \(w\) one-bits, and every signature
fiber has the displayed bound.  Applying that lemma gives
\[
  |P|\le 2^k\Delta^w h^{k-w}.
\]
Replacing \(|P|\) by \(\fwi_k(D^*(G))\) via
\(\noderef{ForwardIndependentTupleCount}\) proves the claim.
\end{proof}
